% !TeX root = Seminararbeit.tex

\documentclass[%
12pt,                % Schriftgröße
paper=a4,            % Papiergröße
captions=tableabove, % Beschriftungen für Tabellen oberhalb
]{scrartcl}

% ----------------------------------------------------
% Essential packages
% ----------------------------------------------------
\usepackage[utf8]{inputenc}
\usepackage[T1]{fontenc}

% ----------------------------------------------------
% Packages for layout adjustments
% ----------------------------------------------------

% Adjust line spacing
\usepackage{setspace}

% Publication quality tables
\usepackage{booktabs}

% ----------------------------------------------------
% Fonts
% ----------------------------------------------------
\usepackage{lmodern}
\renewcommand{\seriesdefault}{m}\selectfont

\newcommand\roboto{\fontfamily{Roboto-LF}\selectfont}
\newcommand*\robotocondensed{\roboto\fontseries{c}\selectfont}

\setkomafont{subject}{\large\robotocondensed}
\addtokomafont{title}{\LARGE}
\addtokomafont{subtitle}{\Large}
\setkomafont{author}{\normalsize\robotocondensed}
\addtokomafont{publishers}{\normalsize\robotocondensed}

% ----------------------------------------------------
% Colors
% ----------------------------------------------------
\usepackage{graphicx}
\usepackage[svgnames]{xcolor}
\definecolor{darkgreen}{rgb}{0.23,0.46,0.23}
\definecolor{smdsblue}{RGB}{0,69,134}

% ----------------------------------------------------
% Internal commands
% ----------------------------------------------------

\usepackage{etoolbox}
\makeatletter
\newcommand{\seminartype}[2]{%
  \subject{%
    Seminar\\
    \textit{\GetTranslationWarn{seminar@#1}}\\
    #2
  }
}
\newcommand{\advisor}[1]{%
  \publishers{%
    \GetTranslation{advisor}: #1\\
    \GetTranslation{institute}}
}
\newcommand{\email}[1]{\gdef\@email{#1}}
\newcommand{\matrno}[1]{\gdef\@matrno{#1}}
\newcommand{\institute}[1]{\gdef\@institute{#1}}
\newcommand{\useAI}[1]{\def\AIUse{#1}}
\makeatother

% Sprachauswahl:
%  main=* setzt die Hauptsprache für das Dokument
%  - ngerman --> deutsch
%  - english --> englisch
\def\languages{main=german,english}
%\def\languages{main=english,ngerman}

% Art und Zeitpunkt des Seminars:
% - SEvS    Software Engineering für verteilte Systeme
% - ML     Machine Learning
% - SEisS    Software Engineering in sicherheitskritischen Systemen
\seminartype{SEvS}{Sommersemester 2024}

% Haupttitel der Arbeit
\title{Recent Approaches to accelerate CP Solving}
% Untertitel der Arbeit -- für Seminararbeiten nicht benötigt
% \subtitle{Concepts, Technologies, and Applications}

% Name, Matrikelnummer und E-Mail-Adresse
\author{---}
\matrno{---}
\email{---e}

% Transparenzangabe zur Verwendung
% künstlicher Intelligenz (KI)-basierter Tools. 
% Mögliche Optionen:

% - No          Keine KI-basierten Tools verwendet:
%               Sämtliche Inhalte sind eigenständig und ohne die
%               Unterstützung von Algorithmen oder Software, 
%               die auf KI basiert, entstanden.

% - Support     Verwendung KI-basierter Tools zur sprachlichen 
%               Verbesserung oder Korrektur des Textes:
%               Dies umfasst die Nutzung von Software zur 
%               Grammatikprüfung, Rechtschreibkorrektur 
%               und stilistischen Optimierung des Textes. 

% - Content     Verwendung KI-basierter Tools zur (teilweisen) 
%               Generierung von Inhalt:
%               Dies beinhaltet die Nutzung von KI für die Erstellung 
%               von Textabschnitten, Konzeptionierung von Ideen oder
%               Bereitstellung struktureller oder inhaltlicher Vorschläge.

\useAI{Content}

% Datum der Abgabe
\date{---}

\advisor{---}

%%% Local Variables:
%%% mode: latex
%%% TeX-master: "Seminararbeit"
%%% End:


% ----------------------------------------------------
% Multi-lingual documents with Babel
% ----------------------------------------------------
\usepackage{csquotes}
\usepackage[\languages]{babel}

% ----------------------------------------------------
% Hyperlinks in PDF documents
% ----------------------------------------------------
\usepackage[%
bookmarks=true,         %
bookmarksopenlevel=1,   %
bookmarksopen=true,     %
bookmarksnumbered=true, %
plainpages=false,       % correct hyperlinks
pdfpagelabels=true,     % view TeX pagenumber in PDF reader
colorlinks=true,        % color highlight links
allcolors=black,        % make all links black by default
urlcolor=smdsblue,      % URL color
]{hyperref}

\makeatletter
\AtEndPreamble{
  \hypersetup{
    pdftitle=\@title,
    pdfauthor=\@author
  }
}
\makeatother

% Provides a solution to the problem with hyperref that links
% to floats actually anchor to the place below the float's caption,
% instead of anchoring to the beginning of the float
\usepackage[all]{hypcap}

% ----------------------------------------------------
% Code listings
% ----------------------------------------------------
\usepackage{listings}
\lstset{%
  frame=single,                             % Add a single line frame around listings
  frameround=ftft,                          % Rounded frame corners on top left and bottom right
  backgroundcolor=\color{gray!5},           % Slight gray shade for listings
  rulecolor=\color{black!30},               % Gray frame outline
  xleftmargin=.125\textwidth,               % Extra left margin
  xrightmargin=.125\textwidth,              % Extra right margin
  basicstyle=\small\ttfamily,               % General font style for listings
  keywordstyle=\bfseries,                   % Font style for keywords
  commentstyle=\color{gray},                % Font style for comments
  stringstyle={},                           % Font style for string literals
  numbers=left,                             % Show line numbers
  stepnumber=1,                             % Step increments for line numbers
  numberstyle={\sffamily\tiny\color{gray}}, % Font style for line numbers
  numbersep=2em,                            % Space between line numbers and code
}

% ----------------------------------------------------
% Bibliography management
% ----------------------------------------------------
\usepackage[%
backend=biber,      % Use biber to process bibliographies
natbib=true,        % Provide natbib-compatible citation commands
sorting=none,       % Sort citations by occurrence in the document
style=numeric-comp, % Use compressed numeric citations, e.g. [1-3; 5]
block=space,        % Add a little spacing inside bibliography entries
]{biblatex}
\addbibresource{literature.bib}

% Use main body font for URLs in bibliography
\urlstyle{same}

% Suppress page numbering on table of contents page(s). Works at least on one-page TOC.
\AtBeginDocument{\addtocontents{toc}{\protect\thispagestyle{empty}}} 

% Intelligent cross-referencing
% Note: Must be loaded at end of preamble (esp. after hyperref)
\usepackage{cleveref}

% ----------------------------------------------------
% Localization / translations
% ----------------------------------------------------
\usepackage{translations}

% Translations for seminar names
\NewTranslation{ngerman}{seminar@SEvS}{Software Engineering für verteilte Systeme}
\NewTranslationFallback{seminar@SEvS}{Software Engineering for Distributed Systems}
\NewTranslation{ngerman}{seminar@MS}{Machine Learning}
\NewTranslationFallback{seminar@MS}{Machine Learning}
\NewTranslation{ngerman}{seminar@SEisS}{Software Engineering in sicherheitskritischen Systemen}
\NewTranslationFallback{seminar@SEisS}{Software Engineering in Safety- and Security-Critical Systems}

% Generic translation used in template
\NewTranslation{ngerman}{advisor}{Betreuer}
\NewTranslation{ngerman}{matrno}{Matrikelnummer}
\NewTranslation{ngerman}{institute}{Softwaremethodik für verteilte Systeme (Prof. Bauer)\\Universität Augsburg}
\NewTranslation{ngerman}{regularlit}{Literatur}
\NewTranslation{ngerman}{onlinelit}{Online-Quellen}
\NewTranslation{ngerman}{honesty@title}{Eidesstattliche Erklärung}
\NewTranslation{ngerman}{honesty@body}{%
  Ich versichere, dass ich die vorliegende Arbeit ohne fremde Hilfe und ohne Benutzung anderer
  als der angegebenen Quellen angefertigt habe, und dass die Arbeit in gleicher oder ähnlicher
  Form noch keiner anderen Prüfungsbehörde vorgelegen hat.\endgraf
  Alle Ausführungen der Arbeit, die wörtlich oder sinngemäß übernommen wurden, sind als solche
  gekennzeichnet.
}
\NewTranslation{ngerman}{aiused@no}{%
  Bei der Erstellung dieses Dokuments wurde keinerlei auf künstliche Intelligenz 
  (KI)-basierte Software verwendet.
}
\NewTranslation{ngerman}{aiused@support}{%
  Bei der Erstellung dieses Dokuments wurde künstliche Intelligenz (KI)-basierte 
  Software ausschließlich zur sprachlichen Verbesserung und Korrektur verwendet. 
}
\NewTranslation{ngerman}{aiused@content}{%
  Bei der Erstellung dieses Dokuments wurde künstliche Intelligenz (KI)-basierte 
  Software zur Generierung von Inhalten verwendet. 
}


% English fallback text
\NewTranslationFallback{advisor}{Advisor}
\NewTranslationFallback{matrno}{Matriculation number}
\NewTranslationFallback{institute}{Software Methodologies for Distributed Systems (Prof. Bauer)\\University of Augsburg}
\NewTranslationFallback{regularlit}{Literature}
\NewTranslationFallback{onlinelit}{Online resources}
\NewTranslationFallback{honesty@title}{Declaration of Academic Honesty}
\NewTranslationFallback{honesty@body}{%
  Hereby, I declare that I have composed the presented paper independently on my own and without
  any other resources than the ones indicated. All thoughts taken directly or indirectly from external
  sources are properly denoted as such.\endgraf
  This paper has neither been previously submitted to another authority nor has it been published yet.
}
\NewTranslationFallback{aiused@no}{%
  No artificial intelligence (AI)-based software was used in the creation of this document.
}
\NewTranslationFallback{aiused@content}{%
  Artificial intelligence (AI)-based software was used for content generation in the 
  creation of this document.
}
\NewTranslationFallback{aiused@support}{%
Artificial intelligence (AI)-based software was used exclusively for linguistic improvement 
and correction in the creation of this document.
}


%%% Local Variables:
%%% mode: latex
%%% TeX-master: "Seminararbeit"
%%% End:


\begin{document}
\pagenumbering{roman}	
% !TeX root = Seminararbeit.tex

\begin{titlepage}
  \onehalfspacing
  \makeatletter
  \vspace*{1em}
  \begin{center}
    \ifdefempty{\@subject}{}{%
      {\usekomafont{subject}\@subject}
      \par\vspace{2em}
    }
    {\usekomafont{title}\@title}
    \ifdefempty{\@subtitle}{}{%
      \par\vspace{.5em}
      {\usekomafont{subtitle}\@subtitle}
    }
    \par\vspace{2em}
    \singlespacing
    {\usekomafont{author}%
      \@author\par
      \GetTranslation{matrno}: \@matrno\par}
    \texttt{\@email}
    \par\vspace{1.5em}
    {\usekomafont{publishers}\@publishers}
  \end{center}
  \makeatother

  \begin{abstract}
    \noindent%
    \paragraph*{\abstractname}
    \textcolor{red}{Bitte Dateien \texttt{settings.tex} und \texttt{titlepage.tex} anpassen.}
    
    Abstract Lorem ipsum dolor sit amet, consectetuer adipiscing elit, sed diam nonummy nibh euismod tincidunt ut laoreet dolore magna aliquam erat volutpat. Ut wisi enim ad minim veniam, quis nostrud exerci tation ullamcorper suscipit lobortis nisl ut aliquip ex ea commodo consequat.
  \end{abstract}

  \vfill
  \centering
  \includegraphics[height=38mm]{figures/uni_siegel}
\end{titlepage}
%%% Local Variables:
%%% mode: latex
%%% TeX-master: "Seminararbeit"
%%% End:


\tableofcontents

\clearpage
\pagenumbering{arabic}


\section{Einleitung: Verwendung von Constraint Programming}
\label{sec:Einleitung: Verwendung von Constraint Programming}
Constraint Programming (CP) spielt in einem weiten Bereichen heutiger modernen
Anwendungen eine große Rolle zur Lösung von Problemen. Dazu gehören unter
anderem: Vehicel Routing, Scheduling, Planung, Konfiguration,
Ressourcenallokation und Kombinatorische Optimierung. Neben dem Klassischen
Problem der Kursplannung \cite{duboi96jo} in der immer schneller werden Zeit,
gib es auch in der Industrie eine Vielzahl von Scheduling Problemen, wie der
Zuweiung von Aufträgen zu Maschinen \cite{gedik16jo}. Auch wurden CP Ansätze für
die Konfiguration von Netzwerken verwendet \cite{ardisjo}. Ein weit Verbreitetes
Kombinatirsches Optimierungsproblem ist das 3-SAT Problem. Es beschreibt die
Frage ob eine Formel in konjunktiver Normalform (KNF) mit Klauseln aus jeweil 3
Literalen erfüllbar ist.  
%\cite{keylist} 
Beispielsweise löst
$x_1=\mathrm{false},x_2=\mathrm{false},x_3=\mathrm{false},x_4=\mathrm{false}$
das folgende 3-SAT Problem: $(\lnot x_1\lor x_2\lor x_3)\land(x_1\lor\lnot
x_2\lor x_4)\land(x_2\lor x_3\lor\lnot x_4)$

Unter dem selben Namen wird Jährlich die MiniZinc Challenge ausgetragen, in der
ein Parkour von Constraint Modelen gelöst werden muss. Am Ende werden die Solver
nach Anzahl der gelösten Modelle, der Zeit und der Qualität der Lösung bewertet.
%\cite{keylist} Vehicel Routing ist in der
heutighen Zeit für viele Unternehmen in der Logistikbranche und dem Flexiblen
Transport von großer Bedeutung. \cite[1]{delec22jo}. Je nach Anwendungen und
Problemstellung können unterschiedliche Constraint Programming Ansätze verwendet
werden. Es exisiteren diveres Toolboxen um Contraint Programming Problem zu
Lösen. Für das Vehicel Routing Problem, welches das Problem beschreibt für n
Fahrzeuge und m Orte den kürzesten gesamt Weg zu finden \cite[222]{labor18joa},
kann beispielseise der IBM ILOG CP Optimzer verwendet werden.\cite{IBMIwe} Je
nach Problemstellung gibe es eine Viehlzahl von weitern Toolboxen
\cite{Solviwea}, Or-tools von Google ist eine weitere zum lösen von Constraint
Programming (CP), Linear Programming (LP), Integer Programming (IP) und Boolean
Satisfiability (SAT) Problemen. \cite{ORToowe}, Geocode ist eine Open Source
Variante basierend auf C++ \cite{GECODwe} und Chuffed is eine Varainte welche
"lazy clause generation" ausnutzt um CP Probleme schneller zu lösen
\cite{Chuff24co}. MiniZinc vereint unter anderem die im Vorherigen beschriebenen
Softwarebibliotheken zu einer Ausdrucksspache zum Lösen von einer Vielzahl von
Lineare Programmierung (LP), Transportprobleme (TP) und SAT Problemen.
\cite{MiniZwe}


%chat gpt 13.06.2024
Transport und Logistik: Zeitplanung und Ressourcenmanagement:
Konfigurationsprobleme: Personaleinsatzplanung:

\cite[1-5]{rossi06bo}

\section{Grundlagen}
\subsection{Constraint Satisfaction Problem und Constraint Optimization Problem} 
\label{sec:Constraint Satisfaction Problem und Constraint Optimization Problem}

finite set of constraints \cite[1]{aptjo} constraint satisfaction problem
\cite[1]{aptjo} (1) problem modeling (2) solve it - has solution - find solution
- find optimal solution
\cite*[postnote]{aptjo}
Constraint Programming (CP), Mixed Integer Linear Programming (MIP) or Boolean
Satisfiability (SAT) solvers



\subsection{Constraint Optimization Problem}
\label{sec: Constraint Optimization Problem}

Neben dem Klassischen Problem der Kursplannung \cite{duboi96jo} in der immer
schneller werden Zeit, gib es auch in der Industrie eine Vielzahl von Scheduling
Problemen, wie der Zuweiung von Aufträgen zu Maschinen \cite{gedik16jo}. Auch
wurden CP Ansätze für die Konfiguration von Netwerken verwendet \cite{ardisjo}.
Ein weit Verbreitetes Kombinatirsches Optimierungsproblem ist das 3-SAT Problem.
Es beschreibt die Frage ob eine Formel in konjunktiver Normalform (KNF) mit
Klauseln aus jeweil 3 Literalen erfüllbar ist.  
%\cite{keylist} 
Beispielsweise löst 
$x_1=\mathrm{false},x_2=\mathrm{false},x_3=\mathrm{false},x_4=\mathrm{false}$
das folgende 3-SAT Problem: 
$(\lnot x_1\lor x_2\lor x_3)\land(x_1\lor\lnot x_2\lor x_4)\land(x_2\lor x_3\lor\lnot x_4)$

Unter dem selben Namen wird Jährlich die MiniZinc Challenge ausgetragen, in der
ein Parkour von Constraint Modelen gelöst werden muss. Am Ende werden die Solver
nach Anzahl der gelösten Modelle, der Zeit und der Qualität der Lösung bewertet.
%\cite{keylist}

\subsection{Boolean Satisfibility problem}
\label{sec: Boolean Satisfibility problem}


\subsection{Mixed Integer Problem}
\label{sec: Mixed Integer Problem}


\section{Recent Approaches To Accelerate Constraint programms}
\label{sec:Recent Approaches To Accelerate Constraint programms}


\subsection{CP and Machine Learning}
\label{sec:CP and Machine Learning}

  
\subsection{Portfolio}
\label{sec:Portfolio}

  
\subsection{Model Based Optization}
\label{sec:Model Based Optization}

  
\subsection{Qunatum Acceleratedt}
\label{sec:Qunatum Accelerated}

  
\subsection{Combine CP and SAT}
\label{sec:Combine CP and SAT}

  
\subsection{Decompostion methods}
\label{sec:Decompostion methods}

  
\subsection{Search Heuristics}
\label{sec:Search Heuristics}


\section{Evaluation Performance Increase}
\label{sec:Evaluation Performance Increase}

MiniZinc vereint unter anderem die im vorherigen beschrieb
Softwarebibliotheken zu einer Ausdrucksspache zum Lösen von einer Viehlzahl von 
LP, Tp und SAT Problemen. \cite{MiniZwe}




% Literaturverzeichnis
\printbibliography[heading=bibintoc]

% Anhang
\include{appendix}

% Eidesstattliche Erklärung
% !TEX root = Seminararbeit.tex

\clearpage
\section*{\GetTranslation{honesty@title}}
\GetTranslation{honesty@body}

\vspace{2em}

\ifdefined\AIUse
    \expandafter\ifstrequal\expandafter{\AIUse}{No}{%
        \GetTranslation{aiused@no}
    }{%
        \expandafter\ifstrequal\expandafter{\AIUse}{Content}{%
            \GetTranslation{aiused@content}
        }{%
            \expandafter\ifstrequal\expandafter{\AIUse}{Support}{%
                \GetTranslation{aiused@support}
            }{% Fallback for undefined values
                \GetTranslation{aiused@no}
            }%
        }%
    }%
\else
    % Fallback, falls \useAI{} nie aufgerufen wurde
    \GetTranslation{aiused@no}
\fi

\vspace{2em}
\makeatletter
Augsburg, \@date
\par\vspace{1.5cm}
(\@author)
\makeatother

%%% Local Variables:
%%% mode: latex
%%% TeX-master: "Seminararbeit"
%%% End:


\end{document}